\section{Fitur Teks Lanjutan: Kutipan, Kode, Subscript, Superscript}

HTML menyediakan elemen khusus untuk menandai berbagai jenis konten khusus seperti kutipan, blok kode, dan notasi matematika. Elemen \code{<blockquote>} digunakan untuk kutipan panjang yang mencakup satu atau lebih paragraf, biasanya ditampilkan browser dengan indentasi di kedua sisi. Atribut opsional \code{cite} pada \code{<blockquote>} dapat berisi URL yang menunjuk ke sumber kutipan tersebut. Untuk kutipan inline yang lebih pendek dan berada dalam aliran teks paragraf, gunakan elemen \code{<q>}. Browser biasanya secara otomatis menambahkan tanda kutip yang sesuai dengan bahasa dokumen (misalnya "..." untuk Inggris, «...» untuk beberapa bahasa Eropa) di sekitar konten \code{<q>}, menghindari kebutuhan menulis tanda kutip secara manual dalam markup \cite{mdnHtmlIntro}.

Teks kode sumber dan output komputer ditandai menggunakan elemen \code{<code>} untuk inline code dan kombinasi \code{<pre>} dengan \code{<code>} untuk blok kode. Elemen \code{<code>} memberikan semantik bahwa konten adalah fragmen kode sumber atau output komputer, yang biasanya ditampilkan dengan font monospaced. Elemen \code{<pre>} (preformatted text) sangat penting untuk blok kode karena mempertahankan spasi, indentasi, dan jeda baris persis seperti dalam source code. Tanpa \code{<pre>}, browser akan mengabaikan spasi berlebih dan menjadikan kode sulit dibaca. Kombinasi \code{<pre><code>...</code></pre>} adalah pola standar untuk menampilkan blok kode dengan formatting terjaga.

Notasi matematika dan kimia sering memerlukan karakter yang diposisikan di atas atau di bawah garis dasar teks. \textbf{Subscript} (karakter rendah) menggunakan elemen \code{<sub>} dan cocok untuk rumus kimia seperti H\textsubscript{2}O atau indeks matematika. \textbf{Superscript} (karakter tinggi) menggunakan elemen \code{<sup>} dan berguna untuk eksponen matematika seperti x\textsuperscript{2}, notasi tanggal (1\textsuperscript{st}), atau catatan kaki. Kedua elemen ini menjaga semantik konten yang memerlukan penurunan atau penaikan posisi karakter, bukan sekadar styling visual.

\textbf{HTML entities} adalah kode khusus untuk menampilkan karakter yang memiliki makna khusus dalam HTML atau karakter yang tidak mudah diketik. Karena karakter \code{<}, \code{>}, dan \code{\char38} memiliki makna khusus dalam HTML (membuka tag, menutup tag, dan memulai entity), mereka harus ditulis sebagai entity \code{\char38lt;}, \code{\char38gt;}, dan \code{\char38amp;} ketika ingin menampilkannya sebagai teks biasa. Entity \code{\char38nbsp;} (non-breaking space) membuat spasi yang tidak akan dipisahkan oleh pemisah baris, berguna untuk mencegah "janda" (kata tunggal di akhir paragraf) atau memisahkan angka dan satuan. Entity \code{\char38copy;} menghasilkan simbol copyright \textcopyright, dan \code{\char38reg;} menghasilkan trademark \textregistered.

\begin{table}[htbp]
\centering
\begin{tabular}{|P{3cm}|P{3cm}|P{5cm}|}
\hline
\textbf{Entity} & \textbf{Hasil} & \textbf{Penggunaan} \\
\hline
\code{\char38lt;} & $<$ & Menampilkan kurung sudut kiri \\
\hline
\code{\char38gt;} & $>$ & Menampilkan kurung sudut kanan \\
\hline
\code{\char38amp;} & \& & Menampilkan ampersand \\
\hline
\code{\char38nbsp;} & Spasi & Non-breaking space \\
\hline
\code{\char38copy;} & \textcopyright & Simbol copyright \\
\hline
\code{\char38reg;} & \textregistered & Simbol registered trademark \\
\hline
\code{\char38hellip;} & \ldots{} & Elipsis horizontal \\
\hline
\code{\char38mdash;} & --- & Em dash (panjang) \\
\hline
\end{tabular}
\caption{HTML Entities Umum}
\label{tab:html-entities}
\end{table}

\begin{webcode}[language=HTML, caption={Contoh Fitur Teks Lanjutan}]
<!-- Kutipan panjang dengan sumber -->
<blockquote cite="https://example.com/article">
  <p>Web adalah medium universal. Semua yang terhubung 
     ke internet di seluruh dunia dapat berkomunikasi 
     tanpa hambatan.</p>
</blockquote>

<!-- Kutipan inline -->
<p>Tim Berners-Lee berkata: 
   <q>The Web does not just connect machines, 
      it connects people</q>.</p>

<!-- Blok kode dengan pre dan code -->
<pre><code>function sapa(nama) {
  return "Halo, " + nama + "!";
}

console.log(sapa("Dunia"));</code></pre>

<!-- Inline code -->
<p>Gunakan fungsi <code>console.log()</code> untuk 
   debugging.</p>

<!-- Rumus kimia dan matematika -->
<p>Rumus air: H<sub>2</sub>O</p>
<p>Rumus kuadrat: ax<sup>2</sup> + bx + c = 0</p>

<!-- Karakter khusus dengan entity -->
<p>5 &lt; 10 dan 15 &gt; 12</p>
<p>&copy; 2026 Perusahaan Contoh &mdash; 
   Hak Cipta Dilindungi</p>
\end{webcode}

